% =====================================================================
%  Installation & Operations Guide: installing, running and maintaining
%  the application, plus every script, setting and common problem.
%  Build: ./scripts/build_docs.sh   (or compile in docs/ with pdflatex)
% =====================================================================
\documentclass[11pt,a4paper]{article}
\usepackage{campusdoc}

\begin{document}

\makedoctitle{Installation \& Operations Guide}{Installing, running and maintaining Campus Events}{CEMP-IG-01}

\tableofcontents
\newpage

% =====================================================================
\section{Overview}

This guide is for whoever installs and runs the Campus Event Management Platform: a lab
assistant, an examiner, or a developer setting up their own copy. It covers system requirements,
installation, day-to-day commands, configuration, the database, and solutions to common problems.

The application is a Flask web app that stores its data in a single SQLite file. Everything is
driven by Bash scripts in \code{scripts/}, which work in \textbf{Git Bash on Windows} and in any
Bash shell on \textbf{Linux or macOS}.

\begin{note}[Quick start]
\begin{lstlisting}[style=shell, backgroundcolor=\color{white}, frame=none]
git clone https://github.com/PrajyotKira/Campus-Event-Management-Platform.git
cd Campus-Event-Management-Platform
./scripts/setup.sh      # venv, packages, database, indexes, admin account
./scripts/seed.sh       # optional: demo data
./scripts/run.sh        # open http://127.0.0.1:5000
\end{lstlisting}
\end{note}

% =====================================================================
\section{System requirements}

\begin{table}[H]
\centering
\small
\begin{tabularx}{\linewidth}{@{}l l X@{}}
\toprule
\textbf{Component} & \textbf{Version} & \textbf{Notes} \\
\midrule
Python      & 3.10 or newer & From \url{https://www.python.org/}. On Windows, tick ``Add Python to PATH''. Developed with Python 3.14; also tested on Python 3.13 (Linux). \\
Bash        & 4 or newer    & Windows: Git Bash, installed with Git for Windows (\url{https://git-scm.com/}). Linux/macOS: built in. \\
Git         & any recent    & To clone the repository. \\
Web browser & any modern    & Chrome, Firefox, Edge or Safari. \\
LaTeX       & optional      & TeX Live or MiKTeX, only to rebuild the PDF documentation. Overleaf also works. \\
\bottomrule
\end{tabularx}
\caption{Requirements.}
\end{table}

SQLite needs no installation: it is built into Python's standard library. The Python packages are
listed in \code{requirements.txt} and installed automatically:

\begin{table}[H]
\centering
\small
\begin{tabular}{@{}l l l@{}}
\toprule
\textbf{Package} & \textbf{Version range} & \textbf{Used for} \\
\midrule
Flask        & \code{>=3.0,<4.0} & web framework, Jinja2 templates, CLI commands \\
Flask-Login  & \code{>=0.6.3}    & login sessions \\
pytest       & \code{>=8.0}      & automated tests \\
Faker        & \code{>=25.0}     & realistic demo data \\
snakeviz     & \code{>=2.2}      & viewing profiler output in the browser \\
\bottomrule
\end{tabular}
\caption{Python dependencies (Werkzeug, Jinja2 and Click come with Flask).}
\end{table}

% =====================================================================
\section{Installation}

\subsection{Get the code}

\begin{lstlisting}[style=shell]
git clone https://github.com/PrajyotKira/Campus-Event-Management-Platform.git
cd Campus-Event-Management-Platform
\end{lstlisting}

You can also download the repository as a ZIP from GitHub and extract it. If the folder name
contains spaces, keep the quotes when you \code{cd} into it.

\subsection{Run the setup script}

\begin{lstlisting}[style=shell]
./scripts/setup.sh
\end{lstlisting}

It performs six steps and prints each one:

\begin{enumerate}
  \item \textbf{Find Python.} Tries \code{python3}, then \code{python}, by actually running
        \code{--version} (on Windows \code{python3} can be a non-working Microsoft Store shortcut).
  \item \textbf{Create the virtual environment} in \code{venv/} (skipped if it already exists).
  \item \textbf{Install requirements} with \code{pip}.
  \item \textbf{Create the database} \code{instance/campus.db} from \code{sql/schema.sql}.
  \item \textbf{Add indexes} from \code{sql/indexes.sql}.
  \item \textbf{Create the admin account} (default \code{admin@campus.local} / \code{admin123}).
\end{enumerate}

\begin{warning}
Running \code{setup.sh} again \textbf{deletes all data}, because step 4 recreates every table.
Back up \code{instance/campus.db} first if you need it (Section~\ref{sec:backup}).
\end{warning}

\subsection{Admin account}
\label{sec:admin}

Choose your own admin login with environment variables:

\begin{lstlisting}[style=shell]
ADMIN_NAME="Event Office" ADMIN_EMAIL=events@iitdh.ac.in ADMIN_PASSWORD='S3cure!pw' ./scripts/setup.sh
\end{lstlisting}

To add another admin later without wiping the database:

\begin{lstlisting}[style=shell]
source venv/Scripts/activate      # Windows Git Bash  (Linux/macOS: source venv/bin/activate)
python -m flask --app app create-admin --name "Second Admin" --email second@campus.local --password secret
\end{lstlisting}

If the email already exists the command does nothing and says so.

\subsection{Load demo data (optional)}

\begin{lstlisting}[style=shell]
./scripts/seed.sh          # 50 events
./scripts/seed.sh 200      # 200 events
\end{lstlisting}

This replaces all venues, events, registrations and student accounts (admins are kept) with:
10 real-sounding campus venues, 100 students (\code{student1@campus.local} to
\code{student100@campus.local}, password \code{student123}), the requested number of events spread
over the next 60 days without any venue clash, and random registrations that never exceed an
event's capacity. The random seed is fixed, so the same command always creates the same data.

\subsection{Start the application}

\begin{lstlisting}[style=shell]
./scripts/run.sh               # http://127.0.0.1:5000
PORT=8000 ./scripts/run.sh     # another port
\end{lstlisting}

Open the printed address in a browser. Press \kbd{Ctrl+C} in the terminal to stop the server. The
server runs in debug mode: it reloads automatically when code changes and shows detailed error
pages, which is right for development and demonstrations.

\subsection{Verify the installation}

\begin{enumerate}
  \item Run \code{./scripts/test.sh}. The last line should read \code{45 passed}.
  \item Open the site, log in as the admin and open \menu{Admin} $\rightarrow$ \menu{Reports}.
  \item Log in as \code{student1@campus.local} in a private window and register for an event.
\end{enumerate}

% =====================================================================
\section{Script reference}

All scripts live in \code{scripts/}, start with \code{set -euo pipefail} (stop on the first error,
treat unset variables as errors, fail a pipeline if any part fails) and can be run from any folder.

\begin{xltabular}{\linewidth}{@{}>{\raggedright\arraybackslash}p{0.31\linewidth} X@{}}
\toprule
\textbf{Command} & \textbf{What it does} \\
\midrule
\endhead
\code{./scripts/setup.sh} & One-time setup (above). Variables: \code{ADMIN_NAME}, \code{ADMIN_EMAIL}, \code{ADMIN_PASSWORD}. \\
\code{./scripts/run.sh} & Activates the venv, checks the database exists, starts Flask in debug mode. Variable: \code{PORT} (default 5000). \\
\code{./scripts/seed.sh [N]} & Replaces non-admin data with demo data with N events (default 50). N must be a positive whole number. \\
\code{./scripts/test.sh [args]} & Runs \code{pytest -v}. Extra arguments go to pytest, e.g.\ \code{./scripts/test.sh -k clash} or \code{./scripts/test.sh tests/test_admin.py}. \\
\code{./scripts/profile.sh <label> [N]} & Builds a separate database \code{profiling/profile.db} with N events (default 2000), profiles the reports and search, saves results in \code{profiling/<label>/}. Variables: \code{RUNS} (timed runs per workload, default 3), \code{APPLY_INDEXES} (1 = add indexes, 0 = skip). \\
\code{./scripts/build_docs.sh} & Compiles every LaTeX document in \code{report/} and \code{docs/} to PDF with \code{latexmk} or \code{pdflatex}. \\
\code{scripts/common.sh} & Not run directly. Shared functions \code{find_python} and \code{activate_venv}, and the variables \code{PROJECT_ROOT} and \code{VENV_DIR}. \\
\bottomrule
\end{xltabular}

% =====================================================================
\section{Configuration}

\subsection{Environment variables}

\begin{table}[H]
\centering
\small
\begin{tabularx}{\linewidth}{@{}l l X@{}}
\toprule
\textbf{Variable} & \textbf{Default} & \textbf{Meaning} \\
\midrule
\code{SECRET_KEY}     & \code{dev-change-me}          & Signs the login cookie. \textbf{Must} be changed for any real deployment. \\
\code{CAMPUS_DB}      & \code{instance/campus.db}     & Path of the SQLite database file. \code{profile.sh} sets it to \code{profiling/profile.db}. \\
\code{PORT}           & \code{5000}                   & Port used by \code{run.sh}. \\
\code{ADMIN_NAME}     & \code{Admin}                  & Admin display name for \code{setup.sh}. \\
\code{ADMIN_EMAIL}    & \code{admin@campus.local}     & Admin login for \code{setup.sh}. \\
\code{ADMIN_PASSWORD} & \code{admin123}               & Admin password for \code{setup.sh}. \\
\code{RUNS}           & \code{3}                      & Timed runs per workload in \code{profile.sh}. \\
\code{APPLY_INDEXES}  & \code{1}                      & Whether \code{profile.sh} adds the indexes. \\
\bottomrule
\end{tabularx}
\caption{Settings read from the environment.}
\end{table}

Set a variable for a single command by writing it in front, e.g.\ \code{PORT=8000 ./scripts/run.sh},
or for the whole terminal session with \code{export SECRET_KEY=...}.

\subsection{Built-in settings}

\code{create_app()} in \code{app/__init__.py} also sets \code{SESSION_COOKIE_SAMESITE="Lax"}, so the
browser does not send the login cookie with form submissions from other websites. The list of event
categories (\emph{Technical, Cultural, Sports, Workshop, Seminar}) is the \code{CATEGORIES} list in
\code{app/events.py}; edit it there to add a category.

% =====================================================================
\section{Flask command-line commands}

The scripts call these commands; you can also run them yourself after activating the virtual
environment. They are always run as \code{python -m flask --app app <command>} because the
\code{flask.exe} launcher can be blocked on Windows.

\begin{table}[H]
\centering
\small
\begin{tabularx}{\linewidth}{@{}l X@{}}
\toprule
\textbf{Command} & \textbf{Effect} \\
\midrule
\code{init-db}       & Drops and recreates all tables from \code{sql/schema.sql}. \textbf{Deletes all data.} \\
\code{add-indexes}   & Creates the indexes in \code{sql/indexes.sql} (safe to repeat) and lists them. \\
\code{create-admin --name N --email E --password P} & Adds an admin account; skips an existing email. \\
\code{seed --events N --students M} & Replaces non-admin data with demo data (defaults 50 and 100). Needs an admin to exist. \\
\code{run --debug --port P} & Starts the development server (what \code{run.sh} does). \\
\code{routes}        & Lists every URL the app answers to (built into Flask). \\
\bottomrule
\end{tabularx}
\caption{Flask CLI commands.}
\end{table}

% =====================================================================
\section{Running without Bash}

If Bash is not available (for example in Windows PowerShell or Command Prompt), run the same steps
by hand from the project folder:

\begin{lstlisting}[style=shell, caption={Manual setup in PowerShell.}]
python -m venv venv
venv\Scripts\Activate.ps1          # Command Prompt: venv\Scripts\activate.bat
python -m pip install -r requirements.txt
python -m flask --app app init-db
python -m flask --app app add-indexes
python -m flask --app app create-admin
python -m flask --app app seed --events 50      # optional demo data
python -m flask --app app run --debug           # http://127.0.0.1:5000
python -m pytest -v                             # tests
\end{lstlisting}

If PowerShell refuses to run \code{Activate.ps1}, run
\code{Set-ExecutionPolicy -Scope CurrentUser RemoteSigned} once.

% =====================================================================
\section{Database operations}

\subsection{Where the data lives}

All data is in one file, \code{instance/campus.db}. It is listed in \code{.gitignore}, so it is never
committed. The \code{instance/} folder is created automatically.

\subsection{Back up and restore}
\label{sec:backup}

Stop the server, then copy the file:

\begin{lstlisting}[style=shell]
cp instance/campus.db "backup-$(date +%Y-%m-%d).db"     # back up
cp backup-2026-11-10.db instance/campus.db              # restore
\end{lstlisting}

\subsection{Reset}

\begin{itemize}
  \item Empty database, keep the venv: \code{python -m flask --app app init-db}, then
        \code{add-indexes} and \code{create-admin}.
  \item Everything from scratch: delete \code{venv/} and \code{instance/}, then run
        \code{./scripts/setup.sh}.
  \item Fresh demo data only: \code{./scripts/seed.sh}.
\end{itemize}

\subsection{Inspect the data}

Any SQLite browser can open the file, for example \emph{DB Browser for SQLite} or the
\code{sqlite3} command-line tool:

\begin{lstlisting}[style=shell]
sqlite3 instance/campus.db
sqlite> .tables
sqlite> SELECT title, start_time FROM events ORDER BY start_time LIMIT 5;
sqlite> .quit
\end{lstlisting}

Do not edit the file while the server is handling requests.

% =====================================================================
\section{Building the documentation}

\begin{lstlisting}[style=shell]
./scripts/build_docs.sh
\end{lstlisting}

This compiles six documents and places each PDF next to its source:

\begin{table}[H]
\centering
\small
\begin{tabular}{@{}l l@{}}
\toprule
\textbf{Source} & \textbf{Document} \\
\midrule
\code{report/report.tex}              & Project report \\
\code{report/slides.tex}              & Presentation slides (Beamer) \\
\code{docs/user_manual.tex}           & User Manual \\
\code{docs/installation_guide.tex}    & Installation \& Operations Guide (this document) \\
\code{docs/developer_guide.tex}       & Developer Guide \\
\code{docs/test_report.tex}           & Test Plan \& Report \\
\bottomrule
\end{tabular}
\caption{Documentation sources.}
\end{table}

The documents in \code{docs/} share the style file \code{docs/campusdoc.sty} and use the screenshots
in \code{report/figures/}. Without a local LaTeX installation, upload the \code{report/} and
\code{docs/} folders together to Overleaf, keeping them side by side, and compile each \code{.tex}
file with pdfLaTeX.

% =====================================================================
\section{Troubleshooting}

\begin{xltabular}{\linewidth}{@{}>{\raggedright\arraybackslash}p{0.36\linewidth} X@{}}
\toprule
\textbf{Problem} & \textbf{Cause and fix} \\
\midrule
\endhead
\code{ERROR: Python 3 was not found.} & Install Python 3.10+ and make sure it is on the PATH. Open a new terminal afterwards. \\
\texttt{\$'\textbackslash r': command not found} when running a script & The script has Windows (CRLF) line endings. The repository's \code{.gitattributes} prevents this on a fresh clone; to repair a copy, run \code{git rm --cached -r scripts && git checkout -- scripts} or convert with \code{dos2unix scripts/*.sh}. \\
\code{Permission denied} running \code{./scripts/x.sh} & Make the scripts executable: \code{chmod +x scripts/*.sh}, or run \code{bash scripts/x.sh}. \\
\code{flask.exe}: Permission denied (Windows) & A security policy blocks the launcher. The scripts already avoid it by using \code{python -m flask}; do the same when typing commands. \\
\code{ERROR: virtual environment not found.} & Run \code{./scripts/setup.sh}. \\
\code{ERROR: database not found.} & Run \code{./scripts/setup.sh} (it creates \code{instance/campus.db}). \\
\code{No admin account found} from \code{seed} & Create one with \code{create-admin}, or run \code{setup.sh}. \\
\code{Address already in use} & Another program uses port 5000 (on macOS often AirPlay Receiver). Use \code{PORT=8000 ./scripts/run.sh}. \\
\code{ModuleNotFoundError: No module named 'flask'} & The virtual environment is not active. Use the scripts, or \code{source venv/Scripts/activate} (Windows) / \code{source venv/bin/activate}. \\
\code{database is locked} & Another program (e.g.\ DB Browser) has the file open with unsaved changes. Close it or write the changes. \\
Logged out after restarting the server & Expected if \code{SECRET_KEY} changed between runs. Set a fixed value. \\
\code{No LaTeX installation found} from \code{build_docs.sh} & Install TeX Live or MiKTeX, or use Overleaf. \\
\bottomrule
\end{xltabular}

% =====================================================================
\section{Deployment notes}

The project is built for development and demonstration. Before exposing it to real users:

\begin{enumerate}
  \item Set a long random \code{SECRET_KEY}, for example the output of\\
        \code{python -c "import secrets; print(secrets.token_hex(32))"}\,.
  \item Do not use \code{--debug}; its interactive error pages must never be public.
  \item Serve the app with a production WSGI server behind HTTPS, for example\\
        \code{waitress-serve --call app:create_app} on Windows, or\\
        \code{gunicorn "app:create_app()"} on Linux.
  \item Change or remove the default admin password and do not load demo data.
  \item Back up \code{instance/campus.db} regularly.
\end{enumerate}

The project report lists further hardening as future work (CSRF tokens, atomic registration for the
last seat).

% =====================================================================
\section{Uninstalling}

Everything lives inside the project folder. Stop the server and delete the folder. No system files,
services or registry entries are created.

\end{document}
